% 04_pipeline.tex — Progressive Reduction Pipeline
\section{Progressive Reduction Pipeline}
\label{sec:pipeline}

The pipeline operates in three stages, progressively compressing the combinatorial decision space while preserving ground-truth coverage (Figure~\ref{fig:pipeline}).

\subsection{Stage 1: Parallel Diagnostic Agents}

Five diagnostic agents analyze the encoded knowledge along complementary analytical dimensions (Table~\ref{tab:agents}).
Each agent operates independently, producing structured evidence packets that identify regions of the decision space warranting further investigation.

\begin{enumerate}
    \item \textbf{Elasticity Agent}: Analyzes own-price and cross-price elasticities to identify sensitive lever-SKU-region combinations.
    \item \textbf{Interaction Agent}: Tests for statistical interaction effects between lever pairs, distinguishing genuine coupling from multi-collinearity artifacts.
    \item \textbf{Constraint Agent}: Maps the feasible region from encoded policy constraints, identifying binding constraints and conflicts.
    \item \textbf{Temporal Agent}: Identifies trends, seasonality, and structural breaks to determine which levers are in transition.
    \item \textbf{Portfolio Agent}: Analyzes SKU-level effects including cannibalization, halo, and substitution patterns.
\end{enumerate}

The parallel agent architecture reduces the input space from \PIPELINEINPUT{} to \PIPELINESTAGE{1} evidence packets while maintaining \GTCOVERAGE{1} ground-truth coverage.

\begin{table}[ht]
\centering
\caption{Diagnostic agent specifications: analytical dimensions and I/O.}
\label{tab:agents}
\small
\begin{tabular}{llp{3.5cm}p{3.5cm}}
\toprule
\textbf{Agent} & \textbf{Dimension} & \textbf{Primary Input} & \textbf{Output} \\
\midrule
Elasticity & Price sensitivity & Own/cross-price elasticities & Lever sensitivity rankings \\
Interaction & Lever coupling & Pairwise interaction statistics & Synergy/conflict flags \\
Constraint & Feasibility & Policy constraint vectors & Feasible region map \\
Temporal & Time dynamics & Trend/seasonality decomposition & Structural break alerts \\
Portfolio & SKU effects & Cross-SKU correlations & Cannibalization/halo map \\
\bottomrule
\end{tabular}
\end{table}

\subsection{Stage 2: Causal Synthesis}

The synthesis layer assembles evidence packets from all five diagnostic agents into candidate interventions.
Evidence is grouped by lever and direction, causal chains are identified (where multiple agents provide corroborating evidence), and conflicting signals are resolved using an epistemic hierarchy: quantitative $>$ policy $>$ expert, unless confidence-weighted evidence overrides.

Each intervention specifies: lever, direction, magnitude, scope, mechanism, evidence trail, and confidence.
Synthesis reduces \PIPELINESTAGE{1} evidence packets to \PIPELINESTAGE{2} structured interventions.

\subsection{Stage 3: Multi-Dimensional Quality Gate}

The quality gate scores each candidate intervention across five dimensions (Table~\ref{tab:quality}) and filters to the top-$K$ recommendations.
Hard constraint violations trigger an automatic veto regardless of other scores.

\begin{table}[ht]
\centering
\caption{Quality gate scoring dimensions.}
\label{tab:quality}
\small
\begin{tabular}{lcp{6cm}}
\toprule
\textbf{Dimension} & \textbf{Weight} & \textbf{Description} \\
\midrule
Evidence Density & 0.25 & Number of independent evidence sources supporting the intervention \\
Constraint Alignment & 0.25 & Compliance with hard/soft policy constraints (veto on hard violation) \\
Actionability & 0.20 & Specificity of lever, direction, magnitude, and scope \\
Testability & 0.15 & Whether the intervention can be A/B tested with measurable outcomes \\
Novelty & 0.15 & Non-obviousness; penalizes findings any single-lever analysis would produce \\
\bottomrule
\end{tabular}
\end{table}

The final output is \PIPELINEOUTPUT{} high-quality interventions with zero hard-constraint violations.

\begin{figure}[ht]
    \centering
    \includegraphics[width=0.9\textwidth]{figures/fig3_pipeline_funnel.pdf}
    \caption{Progressive space reduction: the pipeline compresses ${\sim}10^{18}$ combinations to \PIPELINEOUTPUT{} actionable recommendations while maintaining full ground-truth coverage at every stage.}
    \label{fig:pipeline}
\end{figure}

\begin{figure}[ht]
    \centering
    \includegraphics[width=0.85\textwidth]{figures/fig7_compression_curve.pdf}
    \caption{Compression curve showing candidate space size (log scale) and ground-truth coverage at each pipeline stage.}
    \label{fig:compression}
\end{figure}
